% Source-specific chapter-IV presentation: verified fourth-edition digital reprint.
\chapter{Entrenchment and Presumptive Projectibility}
\label{sec:p15-goodman-entrenchment-and-projectibility}

Goodman's positive proposal uses the history of actual projections to
discriminate between hypotheses which agree on their current positive
instances. The relevant history concerns adoption of hypotheses, not
merely the occurrence of words. It includes unsuccessful projections;
its use is not an inference that earlier success guarantees later success.
The problem is to characterize legitimate projection using this additional
material, rather than demand that the definition mention only the present
evidence and the hypothesis \parencite[IV.1, pp.~84--87]{Goodman1983FourthProjection}.

The observation-history grue construction and the distinction between
sample agreement and support are developed in
\cref{sec:p15-goodman-projectibility}. Here they supply the problem to
which the history-based construction responds.

This presentation follows all six sections of chapter~IV in the fourth
edition. Its title and copyright leaves identify Harvard University Press
and the 1983 edition; the held copy is a digital reprint whose reprint date
is unspecified. The note to the third edition, dated 1973, reports a
substantive revision: the first edition's three rules became two in the
second edition and one in the third. Section~IV.4 was rewritten and IV.5
changed. Goodman credits the joint 1970 report with Robert Schwartz and
Israel Scheffler. The revised construction below is therefore not
attributed as the unmodified first-edition construction
\parencite[pp.~xxi--xxv]{Goodman1983FourthProjection}.

The 1983 note explicitly reaffirms Goodman's conclusion that projectibility
cannot be defined by syntactic or semantic criteria alone. Its positive
alternative adds the history of linguistic practice. This is the source's
philosophical conclusion, not a theorem here excluding every augmented
model which incorporates that history.

\section{History, determined cases and actual adoption}
\label{sec:p15-goodman-projection-history}

\begin{hbdefinition}{def:p15-goodman-time-status}{Time-relative cases and eligibility}
Fix a time $\tau$, a universe of objects or cases, an accepted record of
determinations up to $\tau$, and a domain $\mathcal H$ of formulated
hypotheses. The chapter first treats simple universal hypotheses: a
predicate is ascribed to everything in the universe, or to everything
satisfying an antecedent predicate. For such an $H$, its relevant
instances are divided into
\[
 T_\tau(H),\qquad F_\tau(H),\qquad I_\tau(H),
\]
the cases determined true, determined false, and not yet determined.
These are epistemic classes relative to the record, not a division of
the actual world into true, false and truth-valueless objects.
For a conditional hypothesis, the relevant cases are antecedent cases;
false-antecedent material implications are not counted as favourable
observed cases of the projected consequent.

The hypothesis is \emph{supported} if $T_\tau(H)$ is nonempty,
\emph{violated} if $F_\tau(H)$ is nonempty, and \emph{exhausted} if
$I_\tau(H)$ is empty. Write
\[
 \mathcal E_\tau
 =\{H\in\mathcal H:T_\tau(H)\ne\varnothing,
       \ F_\tau(H)=\varnothing,\ I_\tau(H)\ne\varnothing\}
\]
for the supported, unviolated, unexhausted hypotheses. This symbol is an
editorial abbreviation for the source's standing restriction, not an
additional normative criterion.
\end{hbdefinition}

A hypothesis can be supported and violated simultaneously. A false
hypothesis can remain unviolated until its false instance is determined;
an unviolated hypothesis need not be true. An exhausted true hypothesis
has no remaining cases to which it can now be projected. These distinctions
also prevent treating a later violation or exhaustion as proof that an
earlier projection was illegitimate \parencite[IV.2--3, pp.~89--93]{Goodman1983FourthProjection}.

\begin{hbdefinition}{def:p15-goodman-actual-adoption}{Actual projection is not eligibility alone}
An \emph{actual projection} at a time is an overt formulation and adoption
of the hypothesis as a basis for anticipating the determination of further
cases, while the hypothesis is supported, unviolated and unexhausted.
Goodman does not require assertion of certainty. Mere availability for
projection, even legitimate availability, is different from actual adoption.
An actual projection may subsequently fail.

The record used to assess entrenchment at $\tau$ contains projections
already made by that time. The competitor domain is not restricted to
hypotheses actually projected in that record. In the source's tenseless
use, an \emph{actual hypothesis} is an uttered or inscribed statement;
its occurrence may be past, present or future. That use of ``actual''
must not be confused with ``already projected at $\tau$.''
\end{hbdefinition}

This domain distinction is essential to the expansion side of the theory:
it must allow eligible but unprojected hypotheses to defeat bad projections.
The source's expression \emph{could have been projected} means supported,
unviolated and unexhausted at the specified time; it does not mean that the
projection would have been legitimate. The construction concerns an
assessment relation, not an effective enumeration of every future utterance
\parencite[IV.2, pp.~87--92; IV.4, p.~100 and nn.~9--10]{Goodman1983FourthProjection}.

The account therefore needs both elimination and expansion: actually
adopted hypotheses include bad projections, while eligible and legitimate
hypotheses need not have been adopted. In IV.1 Goodman treats projectibility
as a dispositional counterpart of being projected. He acknowledges that
using projections to define projectibility itself involves projecting
the predicate ``projected.'' His reply is not a universal proof that this
dependency disappears: a specific dispositional problem can be investigated
using an existing stock of accepted information before the general problem
of dispositions is resolved. His image of already spontaneous predictions
being gradually rectified is explanatory framing, not a measured cognitive
mechanism \parencite[pp.~86--87, 92]{Goodman1983FourthProjection}.

Determining a case in the future is not the same as determining a future
event. Historical events may be among $I_\tau(H)$. Likewise acceptance of
an observation statement and acceptance of a projection are separate
inputs: the source sketches adoption rather than supplying a complete
psychological or operational test for it. In some examples a projected
claim about an entire bag is itself accepted as a determined case for an
overhypothesis; inspection of every marble is not stipulated as necessary
in every application \parencite[p.~91; p.~109 n.~19]{Goodman1983FourthProjection}.

\section{Entrenchment attaches to extensions in typed projection roles}
\label{sec:p15-goodman-earned-inherited}

\begin{hbdefinition}{def:p15-goodman-role-entrenchment}{Earned and inherited entrenchment}
The \emph{earned entrenchment} of a predicate is grounded in earlier actual
projections of that predicate and predicates coextensive with it.
Separate its role as an antecedent from its role as a consequent.
For example, repeated projection of ``all $A$ are $B$'' contributes to
$A$ in the antecedent role and to $B$ in the consequent role. Their
entrenchments in the other roles need not be equal.

A \emph{parent} of a predicate $Q$ is a predicate $P$ applying to classes
among which is the extension of $Q$. A parent is thus a class-level
predicate, not merely an arbitrary broader object predicate. For example,
``army division'' is a parent of ``soldier in the 26th division'': the
extension of the latter is one class to which the former applies.
A predicate may have many parents; coextensive predicates have the same
parents.

\emph{Inherited entrenchment} is appraised through the better-entrenched
parents. In the source's coarse comparison it is subsidiary to earned
entrenchment: a marked difference in earned levels is decisive; inherited
levels distinguish predicates whose earned levels are approximately equal.
\end{hbdefinition}

The source initially takes ``much better entrenched'' as a primitive,
coarse comparison. It acts only where the difference is sufficiently
pronounced; it does not impose a numerical scale or decide every close
pair. The later third-edition note also acknowledges importance, variety
and the character of projections as contributors. Thus assigning the
number of literal word occurrences as an exhaustive entrenchment score
would change the theory \parencite[pp.~94--98, 106--107; p.~104 n.~16; p.~xxii]{Goodman1983FourthProjection}.

Coextensive predicate tokens contribute to the record for the same
extension. A newly coined abbreviation can consequently possess the
entrenchment of an established extension. A familiar word whose uses have
different extensions does not accumulate that one extension's history
merely by being repeatedly spoken. Even a name for the word and the
predicate denoted by that word have different extensions and different
records. This is why redefining a ``type'' to include early green tokens
and later blue tokens does not transfer green's entrenchment to grue
\parencite[pp.~94--96, nn.~6--7]{Goodman1983FourthProjection}.

The inheritance example contrasts marbles in a newly found bag $B$ with
marbles in a newly specified haphazard selection, called zig $A$. Both
antecedents can have negligible and equal earned entrenchment. The
bag antecedent nevertheless inherits from the repeatedly projected
classification of bagfuls of marbles. A comparable history for the
parents of the haphazard selection is absent in the stipulated example.
This is a comparison of supplied histories, not a theorem that any
physical container guarantees uniform contents.

\section{Competitors, overriding and three presumptive classifications}
\label{sec:p15-goodman-override}

\begin{hbdefinition}{def:p15-goodman-competitors}{Hypothesis comparison and conflict}
For $H=\forall x(A(x)\to B(x))$ and
$K=\forall x(A'(x)\to B'(x))$, $H$ is better entrenched than $K$
when its antecedent is better entrenched and its consequent no less well
entrenched, or its consequent is better entrenched and its antecedent
no less well entrenched, with the respective projection roles kept fixed.
Write $H\succ_\tau K$ for this supplied comparison. It is not a sum of
the two entrenchments; superiority in one coordinate does not compensate
for inferiority in the other by a stated arithmetic formula.

The source's \emph{conflict} comparison is made among members of
$\mathcal E_\tau$. Neither hypothesis follows from the other together
with the stated facts of their eligibility, and their projections ascribe
to something different predicates such that only one actually applies.
Write $\operatorname{Conf}_\tau(H,K)$ for the symmetric comparison
with these assumptions made explicit. For instance, the grue/green
conflict presupposes an unexamined emerald and an opposed target
ascription. Agreement on the observed objects is not that witness.
\end{hbdefinition}

Conflict is stronger than merely using different words, and it is not
identical with formal inconsistency of two universal sentences in every
structure. Two predicates may both apply to an object, as with green and
round. A hypothesis entailed by another is excluded from the conflict
comparison just specified. These restrictions prevent silently replacing
the source relation with an arbitrary clash of sentence strings
\parencite[pp.~93--94, 99--101]{Goodman1983FourthProjection}.

\begin{hbdefinition}{def:p15-goodman-short-override}{The printed short override condition}
For $H,K\in\mathcal E_\tau$, define the source's printed short
condition by
\[
 \begin{split}
 \operatorname{Ov}_\tau(H,K)\ \Longleftrightarrow\quad&
 \operatorname{Conf}_\tau(H,K)\land H\succ_\tau K\\
 &{}\land\neg\exists L\in\mathcal E_\tau
       (\operatorname{Conf}_\tau(L,H)\land L\succ_\tau H).
 \end{split}
\]
Thus the overrider is not itself confronted by a still better entrenched
eligible conflict. Write $\operatorname{O}_\tau(K)$ if
$\exists H\in\mathcal E_\tau\,\operatorname{Ov}_\tau(H,K)$.
The corresponding \emph{presumptive} classifications are
\[
 \begin{split}
 \operatorname{P}_\tau(H)\ &\Longleftrightarrow\
 H\in\mathcal E_\tau\ \land\
 \forall K\in\mathcal E_\tau
 (\operatorname{Conf}_\tau(H,K)\to\operatorname{O}_\tau(K)),\\
 \operatorname{N}_\tau(H)\ &\Longleftrightarrow\
 H\in\mathcal E_\tau\ \land\neg\operatorname{O}_\tau(H)\ \land\
 \exists K\in\mathcal E_\tau
 (\operatorname{Conf}_\tau(H,K)\land\neg\operatorname{O}_\tau(K)),\\
 \operatorname{U}_\tau(H)\ &\Longleftrightarrow\
 H\notin\mathcal E_\tau\ \lor\operatorname{O}_\tau(H).
 \end{split}
\]
They stand for projectible, nonprojectible and unprojectible in the
source's provisional assessment at a time. Every quantifier over a rival
is over eligible competitors, not only already adopted competitors.
\end{hbdefinition}

\begin{hblemma}{thm:p15-goodman-short-partition}{proved=annex:p15-goodman-short-partition}{Logical consequences of the short definitions}
The displayed short definitions divide $\mathcal H$ into three disjoint
classes. No two members of the projectible class conflict. This result
requires only the symmetric conflict relation and the displayed
definitions; it does not establish the normative adequacy of those
definitions or their suitability for long conflict hierarchies.
\end{hblemma}

The qualifier ``presumptive'' is important. A nonprojectible hypothesis
is not thereby bizarre or illegitimate: two well-entrenched, unresolved
competitors can have that status while an experiment is awaited.
Unprojectibility also does not entail falsity: an unsupported or exhausted
consequence of an accepted true hypothesis is still unprojectible in the
specified sense. Source p.~104 distinguishes a normal capacity to derive
credibility from support from projectibility under the actual evidence;
the latter is the principal sense of the chapter. An adverse record can
remove it from a normally usable hypothesis, and a sufficiently favourable
record can confer it on an ordinarily suspect one.

\begin{proofboundary}{def:p15-goodman-hierarchy-scope}{The short condition and the longer-hierarchy variant}
At p.~101 n.~11 Goodman explicitly qualifies the short override condition
to hierarchies of at most three successively better-entrenched conflicting
hypotheses. For longer hierarchies he gives a different extension: a
hypothesis is overridden when it is the bottom member of a hierarchy
having an even number of members, maximal because it cannot be extended
upward, and minimal because each member conflicts only with adjacent
members. The parity condition, maximality and absence of nonadjacent
conflicts must all be retained. The short formula must not be used as
if it supplied that longer-hierarchy extension.
\end{proofboundary}

\begin{hblemma}{thm:p15-goodman-parity-hierarchy}{proved=annex:p15-goodman-parity-hierarchy}{The distinction on an isolated finite hierarchy}
Suppose the entire eligible conflict component is
$H_0,\ldots,H_{n-1}$, $n\ge1$, with successively increasing entrenchment,
conflicts exactly between neighbouring members and no external conflicts.
Under the source's parity extension, $H_i$ is overridden exactly when
$n-i$ is even. Applying the same classification clauses with this
overridden predicate gives alternating unprojectible and projectible
members, with the top member projectible and no nonprojectible members.
For the four-member path, in contrast, the printed short condition makes
only $H_2$ overridden, $H_3$ projectible, and $H_0,H_1$ nonprojectible.
\end{hblemma}

The isolated component hypothesis makes the comparison well defined.
The result is not an unmentioned resolution rule for arbitrary branching
graphs or infinite hierarchies. The source's argument about a small number
of markedly distinct entrenchment levels motivates the qualification;
it is not a proof that every competitor domain is finite.

\section{The revised source's cases}
\label{sec:p15-goodman-revised-cases}

The letters $H_1,\ldots,H_8,K$ below retain the source's hypothesis
names. They are not entrenchment values. Let $t$ be the examples' fixed
cutoff. The comparisons assume the source's marked entrenchment differences,
the stated evidence and target-conflict witnesses, and no further
undefeated relevant rivals. Adding other eligible competitors can change
the classifications. These stipulated comparisons are not measurements
of contemporary scientific practice.

\begin{hbexample}{def:p15-goodman-h1-h2}{An unadopted rival and the green--shape variants}
Put $H_1$: all emeralds are grue, and $K$: all emeralds are green.
Given positive green observations before $t$ and remaining unexamined
emeralds, both are eligible. Their antecedent is the same and green is
stipulated much better entrenched than grue as a consequent. An eligible
$K$ can override $H_1$ even if $K$ was not actually adopted. Earlier
actual projections establish the predicate comparison; current eligibility
establishes the competitor's availability.

Next let \emph{grund} mean green for things examined by the cutoff and
round for the remaining things. Put $H_2$: all emeralds are grund.
Green and round can both apply, so $K$ need not conflict with $H_2$.
The cases separate as follows.
\begin{enumerate}
\item If the examined emeralds are also determined square, $H_3$, all
emeralds are square, is eligible. With the specified opposed square/round
target ascription and superior entrenchment of square, it overrides $H_2$.
\item If shape is unexamined, $H_3$ is unsupported. If some examined
emeralds are square and others not square, it is violated. Neither status
permits $H_3$ to override $H_2$.
\item Let \emph{grare} mean green for the early examined things and square
for the others. $H_4$, all emeralds are grare, has the same early colour
support as $H_2$. With their approximately equal entrenchment and no
competent overrider, the opposed target shapes make both nonprojectible.
The same status can characterize well-entrenched scientific competitors
awaiting an experiment; it is not by itself a verdict of illegitimacy.
\item If all examined emeralds are instead green and round, $H_3$ is
violated and $H_5$, all emeralds are round, is eligible. $H_5$ overrides
square-target rivals such as $H_4$, and $H_2$ can be projectible. The
conjunction of $K$ and $H_5$ entails $H_2$: an early case meets its green
branch and a remaining case its round branch.
\end{enumerate}
In the mixed-shape case one can also retreat to the weaker claim that
emeralds are square or round. The stronger alternatives can remain
nonprojectible while this supported, unviolated weaker claim avoids
their conflict. A statistical hypothesis about shape distribution could
produce a different comparison, but it requires further definitions of
support, violation and conflict not supplied by the universal construction.
\parencite[IV.4, pp.~99--104 and n.~13]{Goodman1983FourthProjection}
\end{hbexample}

\begin{hbexample}{def:p15-goodman-emeruby-cases}{Antecedent changes and a defeated overrider}
Let \emph{emeruby} apply to emeralds examined for colour before $t$ and
rubies not examined for colour before $t$. Colour examination and shape
determination must be distinguished. Put $H_6$: all emerubies are green.
Early green emeralds support it without a direct ruby observation.

If some colour-examined rubies have all been determined red, $H_7$, all
rubies are red, is eligible. Its better-entrenched antecedent, with the
consequent comparison stipulated, allows it to override $H_6$. If no
such ruby observations are available, other comparisons remain possible:
blue sapphire observations support the equally contrived blue
\emph{sapphiruby} alternative; a black object such as the Eiffel tower
supports the black \emph{Eiffelruby} alternative. Such comparisons can
leave $H_6$ nonprojectible without a better-entrenched direct ruby rival.

The source also stipulates the extreme background in which the only
colour determinations anywhere are of green emeralds. Then even the
broader claim that everything is green meets the favourable conditions
of its example, and entails $H_6$. This illustrates dependence on the
\emph{total stated} record, not a licence to omit known nongreen objects.

Finally $H_8$, all emerubies are \emph{gred}, combines the changed
antecedent with a consequent taking the green branch on early
colour-examined things and the red branch on the others. With only the
early green emerald determinations, $H_6$ overrides $H_8$. Adding the
red ruby determinations makes $H_7$ eligible; $H_7$ now overrides $H_6$,
so the former overrider no longer blocks $H_8$. The conjunction of $K$
and $H_7$ entails $H_8$ by its two branches. The analogous claim that all
emerubies are grund is not projectible if neither emerald nor ruby shapes
are determined, or if some are determined not round. Under the source's
stipulated green emerald and round ruby record, it can instead be projectible.
\parencite[IV.4, pp.~104--106]{Goodman1983FourthProjection}
\end{hbexample}

These cases do not establish a general closure principle making every
consequence of a projectible hypothesis projectible. A consequence can be
unsupported or exhausted. Nor is normal projectibility identical with
classification under every record. The changing records explain why a
suspect predicate is not rejected solely for its spelling or novelty.

\section{Comparative projectibility and positive overhypotheses}
\label{sec:p15-goodman-comparative}

\begin{hbdefinition}{def:p15-goodman-positive-overhypothesis}{The class-level parent relation in an overhypothesis}
A positive overhypothesis of $H$ has antecedent and consequent predicates
which are parents of the antecedent and consequent predicates of $H$,
respectively. For a bag $B$ in a stack $S$, consider $H_9$: every marble
in $B$ is red, and $G$: every bagful in $S$ is uniform in colour. The
consequent of $G$ applies to classes, including the extension of ``red
marble''; its antecedent applies to bagful classes, including the marbles
in $B$.
\end{hbdefinition}

Assume $H_9$ is already presumptively projectible. Its newly specified
bag antecedent has little earned and only moderate inherited entrenchment,
so its initial degree of projectibility need not be high. Uniformity
observations of other bags in $S$, possibly in different colours, supply
positive cases for $G$. The intermediate projected case is that $B$ is
uniform in colour. It reinforces extending the particular red observations
within $B$ to the remaining marbles. The other bags' uniformity does not
choose red rather than blue for $B$; the red cases within $B$ do that work.
Positive cases and projectibility are distinct contributors to confirmation
\parencite[IV.5, pp.~108--110]{Goodman1983FourthProjection}.

Other positive overhypotheses include uniformity for small bags, uniformity
for bags in Utah, or the claim that each bagful in $S$ is uniformly of
some warm colour, provided the parent relations hold. Logical consequence
does not automatically transmit the full degree
of an overhypothesis's projectibility. The overhypothesis must itself be
presumptively projectible and appreciably projectible to reinforce the
target substantially.

The counterexample is \emph{bagleet}, applying to naval fleets and the
single bag $B$. Uniform-colour fleet observations support ``every bagleet
is uniform'' without directly observing $B$. The contrived \emph{bagmark}
classification applies to makes of cars and $B$; mixed-colour car observations
support an opposed claim. Many favourable fleet cases do not by themselves
make the bagleet overhypothesis projectible. Ignoring its rival recreates
the original difficulty at the class level
\parencite[pp.~110--112]{Goodman1983FourthProjection}.

The source's further overhypothesis $J$ says that each stack of marbles
in Utah is homogeneous in colour variegation: each stack's bags are all
uniform or all mixed. Goodman argues that the poorly entrenched parent
predicates above this variegation predicate leave higher overhypotheses
without appreciable initial projectibility. For a zig-$A$ target with
negligible earned and inherited entrenchment, the same objection propagates
upward: a positive overhypothesis must contain a parent of the target
antecedent, but no appreciable independent source of reinforcement is
reached. This is the source's argument that appraisal need not continue
indefinitely \parencite[pp.~112--114]{Goodman1983FourthProjection}.

\begin{hblemma}{thm:p15-goodman-regress-boundary}{proved=annex:p15-goodman-regress-boundary}{A precise boundary for the regress argument}
Consider an independently specified dependency hierarchy, with an upward
relation and no node initially marked appreciable. Suppose a node can be
finally appreciable only if some node directly above it is finally
appreciable. If upward paths from each node have a finite bound on their
length, no node is finally appreciable. Without the bounded-path condition
the conclusion need not follow: on the infinite chain of natural numbers,
with $i+1$ directly above $i$, marking every node finally appreciable
satisfies the stated necessary condition.
\end{hblemma}

This conditional result identifies a boundary, not a normative criterion.
The source supplies no numerical reinforcement function, least-fixed-point
convention for arbitrary infinite hierarchies, or termination proof for an
appraisal algorithm. Its negligible-index argument is not a mathematical
theorem of that wider scope. A modest but nonnegligible initial index does
not impose an upper bound: sufficiently projectible overhypotheses can
increase it substantially. Further projections and the violation of
competitors can also change the assessment. Novelty is not permanent exclusion.

\section{Specificity, negative overhypotheses and unresolved comparisons}
\label{sec:p15-goodman-negative-overhypotheses}

Three separate contributors to an effective overhypothesis are its degree
of projectibility, its supporting evidence and its specificity. Holding
two fixed, effectiveness varies with the third. Observations of bags from
$S$ support both $G$, uniformity within $S$, and a broader Utah-wide
uniformity claim. They have a more specific bearing on a new bag in $S$
than on a bag in another stack $T$. More general information can matter
more when it supplies substantially more evidence, including cases from
$T$. Overlapping antecedent classes also require avoiding duplicated
evidence \parencite[pp.~114--116 and n.~21]{Goodman1983FourthProjection}.

These are ceteris-paribus comparisons, not an equation adding three
measured quantities. The source does not define a unified numerical
measure of overall projectibility. No two presumptively projectible
hypotheses conflict. Presumptively nonprojectible competitors can instead
have similar initial entrenchment indices but receive greatly different
reinforcement. One can then \emph{outweigh} the other, a comparison distinct
from overriding through markedly greater entrenchment. Ties can still
require an experiment \parencite[pp.~117--118 and n.~22]{Goodman1983FourthProjection}.

\begin{hbdefinition}{def:p15-goodman-negative-overhypothesis}{Negative need not mean syntactically negated}
A negative overhypothesis has a parent of the target antecedent, but its
consequent is complementary to a parent of the target consequent in the
relevant class domain. For $H_9$, ``every bagful in $S$ is mixed in colour''
is negative: mixed applies to precisely the relevant bagfuls which are
not uniform. Its relation to the target is negative despite the absence
of an overt negation sign. Effective negative overhypotheses weaken the
target's projectibility, subject to the same projectibility, evidence
and specificity considerations.
\end{hbdefinition}

Given a nonempty determined bag case in the example, support for universal
mixedness violates universal uniformity, and conversely. They cannot both
be supported, unviolated effective universal overhypotheses at once. If
some bags are uniform and others mixed, neither universal adequately
represents the record. A statistical overhypothesis about proportions
could transmit such mixed evidence, but its definitions require a further
theory which Goodman expressly leaves undeveloped
\parencite[pp.~118--120]{Goodman1983FourthProjection}.

A high initial entrenchment index therefore does not guarantee a high
final degree: adverse correlative information can weaken it. Negative
overhypotheses also let equally entrenched nonprojectible rivals receive
different final assessments. The chapter contemplates extending comparisons
to hypotheses overridden at the presumptive stage but does not pursue it.
It explicitly distinguishes its definitional account from describing or
prescribing a practical procedure; it asserts no mechanical search,
completeness theorem or resource bound \parencite[p.~120]{Goodman1983FourthProjection}.

\section{Closing proposals and their argumentative status}
\label{sec:p15-goodman-survey}

The positive account locates inductive validity partly in linguistic
practice: which extensions have actually been anticipated and how those
anticipations are organized. Its defence separates two objections. It
denies that entrenchment must be explained by antecedently recognized,
language-independent projectibility. It also denies that the entrenchment
account guarantees true future projections. The former is Goodman's
philosophical position; the latter is an explicit limitation, not an
exception to a proved reliability theorem
\parencite[pp.~96--99, 120--121]{Goodman1983FourthProjection}.

Goodman compares the recurrence of explicit term-projection with Hume's
recurrence of experienced features. His comparison with Kant concerns
the organization of what is presented, but he locates that organization
in language use rather than in an inevitable, immutable cognitive form.
These are philosophical comparisons, not empirical demonstrations of
either historical theory \parencite[pp.~96--97]{Goodman1983FourthProjection}.

The counterfactual proposal concerns a copper penny and the coins actually
in a pocket on VE day. Suppose the actual pocket coins were silver. Even
making that pocket generalization lawlike through further uniform-material
pocket observations does not, Goodman argues, make it correct that putting
the copper penny there would have made it silver. The opposed conditional
that the penny would remain copper is sustained by the stronger
material-persistence generalization. The example motivates comparing
conditionals and their conflicts, rather than saying that any true law
sustains every associated counterfactual. Goodman calls the opposed
copper-preserving conditional a \emph{semifactual}: its consequent retains
the actual coin's copper material. The match example suggests a similar
conflict-based explanation for rejecting the claim that, had that match
been scratched, it would not have been dry. The chapter does not supply
a complete match background or a complete counterfactual semantics at
this point. These are proposed applications of comparing conditionals,
not independently proved adequacy results
\parencite[pp.~121--122 and n.~23]{Goodman1983FourthProjection}.

For kinds, the closing suggestion relates greater kinship to a more
specific, better-entrenched predicate applying to the things compared.
For randomness it contrasts bags drawn only from a stack's top layer with
bags drawn from many layers and positions. A narrower, entrenched top-layer
classification applies to every member of the first sample. The comparison
stipulates that an equally narrow classification covering the more diverse
second sample has less inherited entrenchment. This suggested
sense of nonrandom likeness is not a definition of probabilistic random
sampling, independence or exchangeability. Simplicity is another proposed
application, not a supplied complexity measure
\parencite[pp.~122--124]{Goodman1983FourthProjection}.

\begin{proofboundary}{def:p15-goodman-entrenchment-boundary}{Philosophical defence and formal consequences}
The revised source supplies the time-relative statuses, competitor
comparisons, short override condition, qualified longer-hierarchy variant,
case analyses and comparative resources presented here. Its primitive
entrenchment comparisons, projection record and conflict assumptions are
premises of an application. Their normative authority is argued
philosophically, not proved by the three mathematical lemmas. Statistical
extensions, a unified numerical overall measure, arbitrary branching or
infinite hierarchy procedures and the closing counterfactual, kind,
randomness and simplicity proposals are not completed theories in this
chapter. The original first- and second-edition rules remain distinct
from this revised text.
\end{proofboundary}
